% 01_soc_overview.tex -- the whole SoC on one page; every numbered block has its own figure.
% Build: tectonic 01_soc_overview.tex
\documentclass[tikz,border=10pt,12pt]{standalone}
% soc_style.tex -- shared look for the SoC figures (same palette as the slide decks and the
% core register placement figure). Every figure does % soc_style.tex -- shared look for the SoC figures (same palette as the slide decks and the
% core register placement figure). Every figure does % soc_style.tex -- shared look for the SoC figures (same palette as the slide decks and the
% core register placement figure). Every figure does \input{soc_style}.
\usepackage[T1]{fontenc}
\usepackage[scaled=0.95]{helvet}
\renewcommand{\familydefault}{\sfdefault}
\usepackage{booktabs}
\usepackage{amsmath}
\usepackage{amssymb}
\usetikzlibrary{arrows.meta,calc,positioning,fit,shapes.geometric,shapes.misc}

% ---------------- palette ----------------
\definecolor{pcb}{HTML}{0B2E33}
\definecolor{pcbmid}{HTML}{1F5C61}
\definecolor{tint}{HTML}{E6F0EF}
\definecolor{copper}{HTML}{B5651D}
\definecolor{coppert}{HTML}{F6E8DC}
\definecolor{ink}{HTML}{1F2A2E}
\definecolor{muted}{HTML}{5B6B70}
\definecolor{hardf}{HTML}{DCE6F2}
\definecolor{harde}{HTML}{2F4A7A}
\definecolor{good}{HTML}{2E7D32}
\definecolor{bad}{HTML}{B3261E}
% one colour per kind of block: border (d) / fill (f)
\definecolor{cpud}{HTML}{2B5C9E} \definecolor{cpuf}{HTML}{D6E4F5}
\definecolor{gpud}{HTML}{2E7D32} \definecolor{gpuf}{HTML}{D5EDD0}
\definecolor{npud}{HTML}{0B2E33} \definecolor{npuf}{HTML}{BFE0DB}
\definecolor{memd}{HTML}{5E3A8C} \definecolor{memf}{HTML}{E4DAF3}
\definecolor{busd}{HTML}{B5651D} \definecolor{busf}{HTML}{F6E8DC}
\definecolor{prd}{HTML}{546E7A}  \definecolor{prf}{HTML}{E3EAED}
\definecolor{secd}{HTML}{B3261E} \definecolor{secf}{HTML}{F7D9D6}
\definecolor{pwrd}{HTML}{8A6800} \definecolor{pwrf}{HTML}{FFF0B8}
\definecolor{medd}{HTML}{8E2463} \definecolor{medf}{HTML}{F5D6E6}
\definecolor{skd}{HTML}{2E7D32}  \definecolor{skf}{HTML}{BFE3B4}   % skew buffers, as in the register figure

\tikzset{
  blk/.style 2 args={draw=#1, fill=#2, rounded corners=4pt, thick, align=center,
                     font=\scriptsize, inner sep=3pt},
  cpu/.style={blk={cpud}{cpuf}},  gpu/.style={blk={gpud}{gpuf}},
  npu/.style={blk={npud}{npuf}},  mem/.style={blk={memd}{memf}},
  bus/.style={blk={busd}{busf}},  per/.style={blk={prd}{prf}},
  sec/.style={blk={secd}{secf}},  pwr/.style={blk={pwrd}{pwrf}},
  med/.style={blk={medd}{medf}},
  comb/.style ={draw=copper, fill=coppert, rounded corners=2pt, thick, align=center,
                font=\scriptsize, inner sep=2pt},                       % combinational logic
  hard/.style ={draw=harde, fill=hardf, thick, dashed, align=center,
                font=\scriptsize, inner sep=2pt},                       % hard macro / off-chip
  sram/.style ={draw=memd, fill=memf, thick, align=center, font=\scriptsize,
                inner sep=2pt},                                         % SRAM array, not flip-flops
  ff/.style   ={draw=#1, fill=#1!30, line width=0.35pt},                % one flip-flop square
  ffgone/.style={draw=muted, fill=white, densely dashed, line width=0.35pt},
  panel/.style={draw=muted, rounded corners=6pt, line width=0.9pt},
  ptitle/.style={anchor=north west, font=\normalsize\bfseries, text=ink, inner sep=0pt},
  note/.style ={font=\scriptsize, text=ink, align=left, inner sep=0pt},
  lab/.style  ={font=\tiny, text=muted, inner sep=1pt},
  axi/.style  ={{Stealth[length=2.6mm]}-{Stealth[length=2.6mm]}, line width=2.2pt, copper},
  axi1/.style ={-{Stealth[length=2.6mm]}, line width=2.2pt, copper},
  apb/.style  ={{Stealth[length=1.8mm]}-{Stealth[length=1.8mm]}, line width=1pt, copper},
  apb1/.style ={-{Stealth[length=1.8mm]}, line width=1pt, copper},
  dat/.style  ={-{Stealth[length=1.8mm]}, thick, pcbmid},
  ctl/.style  ={-{Stealth[length=1.8mm]}, densely dashed, thick, muted},
  irq/.style  ={-{Stealth[length=2mm]}, densely dotted, line width=1.3pt, bad},
  num/.style  ={circle, fill=copper, text=white, font=\scriptsize\bfseries, inner sep=1pt,
                minimum size=0.46cm},
  ours/.style ={draw=bad, line width=1.2pt, densely dashed, rounded corners=4pt},
}

% ---------------- helpers ----------------
\newcommand{\rn}[1]{\texttt{#1}}
% \ffrow{x}{y}{n}{cell}{style}: n flip-flops in a row, bottom-left corner at (x,y)
\newcommand{\ffrow}[5]{\foreach \q in {1,...,#3}{\draw[#5] ({#1+(\q-1)*#4},{#2}) rectangle ++(#4,#4);}}
% \ffcol{x}{y}{n}{cell}{style}: n flip-flops in a column, top-left corner at (x,y)
\newcommand{\ffcol}[5]{\foreach \q in {1,...,#3}{\draw[#5] ({#1},{#2-\q*#4}) rectangle ++(#4,#4);}}
% \ffgrid{x}{ytop}{n}{per row}{cell}{style}: n flip-flops, filled row by row from the top-left
\newcommand{\ffgrid}[6]{\pgfmathtruncatemacro{\gl}{#3-1}%
  \foreach \b in {0,...,\gl}{\pgfmathtruncatemacro{\gc}{mod(\b,#4)}\pgfmathtruncatemacro{\gr}{div(\b,#4)}%
    \draw[#6] ({#1+\gc*#5},{#2-(\gr+1)*#5}) rectangle ++(#5,#5);}}

% every figure: same page size, title, footer
\newcommand{\canvas}{\useasboundingbox (-0.3,-0.2) rectangle (38.6,30.4);}
\newcommand{\figtitle}[3]{%
  \node[num, minimum size=0.9cm, font=\Large\bfseries] at (0.45,29.55) {#1};
  \node[anchor=west, font=\LARGE\bfseries, text=pcb] at (1.15,29.55) {#2};
  \node[anchor=west, font=\small, text=ink] at (0,28.7) {#3};}
\newcommand{\figfoot}[1]{\node[anchor=west, font=\scriptsize, text=muted] at (0,0.25) {#1};}
% \panel{x0}{y0}{x1}{y1}{title}
\newcommand{\panel}[5]{\draw[panel] (#1,#2) rectangle (#3,#4);
  \node[ptitle] at ({#1+0.3},{#4-0.25}) {#5};}

% ---------------- timing diagrams: 1 cycle = \cw cm, low at y, high at y+0.45 ----------------
\newcommand{\cw}{1.0}
\newcommand{\wclk}[3]{% x0 y ncycles
  \foreach \k in {1,...,#3}{\draw[pcb, thick] ({#1+(\k-1)*\cw},#2) -- ++(0,0.45) -- ++(0.5*\cw,0)
     -- ++(0,-0.45) -- ++(0.5*\cw,0);}}
\newcommand{\wsig}[3]{% x0 y {level per cycle, e.g. 0,1,1,0}
  \gdef\wprev{x}%
  \foreach \v [count=\k from 0] in {#3}{%
    \draw[pcb, thick] ({#1+\k*\cw},{#2+0.45*\v}) -- ++(\cw,0);
    \if x\wprev\else\ifnum\wprev=\v\else\draw[pcb, thick] ({#1+\k*\cw},#2) -- ++(0,0.45);\fi\fi
    \global\let\wprev\v}}
\newcommand{\wbus}[5]{% x0 y c0 c1 text: one value held from cycle c0 to cycle c1
  \draw[pcbmid, thick, fill=white] ({#1+(#3)*\cw},{#2+0.225}) -- ({#1+(#3)*\cw+0.08},{#2+0.45})
    -- ({#1+(#4)*\cw-0.08},{#2+0.45}) -- ({#1+(#4)*\cw},{#2+0.225}) -- ({#1+(#4)*\cw-0.08},#2)
    -- ({#1+(#3)*\cw+0.08},#2) -- cycle;
  \node[font=\tiny\ttfamily] at ({#1+((#3)+(#4))*\cw/2},{#2+0.225}) {#5};}
\newcommand{\widle}[4]{% x0 y c0 c1: bus idle (no value)
  \draw[muted, thick] ({#1+(#3)*\cw},{#2+0.225}) -- ({#1+(#4)*\cw},{#2+0.225});}
\newcommand{\wname}[3]{\node[anchor=east, font=\scriptsize\ttfamily] at ({#1-0.15},{#2+0.225}) {#3};}
.
\usepackage[T1]{fontenc}
\usepackage[scaled=0.95]{helvet}
\renewcommand{\familydefault}{\sfdefault}
\usepackage{booktabs}
\usepackage{amsmath}
\usepackage{amssymb}
\usetikzlibrary{arrows.meta,calc,positioning,fit,shapes.geometric,shapes.misc}

% ---------------- palette ----------------
\definecolor{pcb}{HTML}{0B2E33}
\definecolor{pcbmid}{HTML}{1F5C61}
\definecolor{tint}{HTML}{E6F0EF}
\definecolor{copper}{HTML}{B5651D}
\definecolor{coppert}{HTML}{F6E8DC}
\definecolor{ink}{HTML}{1F2A2E}
\definecolor{muted}{HTML}{5B6B70}
\definecolor{hardf}{HTML}{DCE6F2}
\definecolor{harde}{HTML}{2F4A7A}
\definecolor{good}{HTML}{2E7D32}
\definecolor{bad}{HTML}{B3261E}
% one colour per kind of block: border (d) / fill (f)
\definecolor{cpud}{HTML}{2B5C9E} \definecolor{cpuf}{HTML}{D6E4F5}
\definecolor{gpud}{HTML}{2E7D32} \definecolor{gpuf}{HTML}{D5EDD0}
\definecolor{npud}{HTML}{0B2E33} \definecolor{npuf}{HTML}{BFE0DB}
\definecolor{memd}{HTML}{5E3A8C} \definecolor{memf}{HTML}{E4DAF3}
\definecolor{busd}{HTML}{B5651D} \definecolor{busf}{HTML}{F6E8DC}
\definecolor{prd}{HTML}{546E7A}  \definecolor{prf}{HTML}{E3EAED}
\definecolor{secd}{HTML}{B3261E} \definecolor{secf}{HTML}{F7D9D6}
\definecolor{pwrd}{HTML}{8A6800} \definecolor{pwrf}{HTML}{FFF0B8}
\definecolor{medd}{HTML}{8E2463} \definecolor{medf}{HTML}{F5D6E6}
\definecolor{skd}{HTML}{2E7D32}  \definecolor{skf}{HTML}{BFE3B4}   % skew buffers, as in the register figure

\tikzset{
  blk/.style 2 args={draw=#1, fill=#2, rounded corners=4pt, thick, align=center,
                     font=\scriptsize, inner sep=3pt},
  cpu/.style={blk={cpud}{cpuf}},  gpu/.style={blk={gpud}{gpuf}},
  npu/.style={blk={npud}{npuf}},  mem/.style={blk={memd}{memf}},
  bus/.style={blk={busd}{busf}},  per/.style={blk={prd}{prf}},
  sec/.style={blk={secd}{secf}},  pwr/.style={blk={pwrd}{pwrf}},
  med/.style={blk={medd}{medf}},
  comb/.style ={draw=copper, fill=coppert, rounded corners=2pt, thick, align=center,
                font=\scriptsize, inner sep=2pt},                       % combinational logic
  hard/.style ={draw=harde, fill=hardf, thick, dashed, align=center,
                font=\scriptsize, inner sep=2pt},                       % hard macro / off-chip
  sram/.style ={draw=memd, fill=memf, thick, align=center, font=\scriptsize,
                inner sep=2pt},                                         % SRAM array, not flip-flops
  ff/.style   ={draw=#1, fill=#1!30, line width=0.35pt},                % one flip-flop square
  ffgone/.style={draw=muted, fill=white, densely dashed, line width=0.35pt},
  panel/.style={draw=muted, rounded corners=6pt, line width=0.9pt},
  ptitle/.style={anchor=north west, font=\normalsize\bfseries, text=ink, inner sep=0pt},
  note/.style ={font=\scriptsize, text=ink, align=left, inner sep=0pt},
  lab/.style  ={font=\tiny, text=muted, inner sep=1pt},
  axi/.style  ={{Stealth[length=2.6mm]}-{Stealth[length=2.6mm]}, line width=2.2pt, copper},
  axi1/.style ={-{Stealth[length=2.6mm]}, line width=2.2pt, copper},
  apb/.style  ={{Stealth[length=1.8mm]}-{Stealth[length=1.8mm]}, line width=1pt, copper},
  apb1/.style ={-{Stealth[length=1.8mm]}, line width=1pt, copper},
  dat/.style  ={-{Stealth[length=1.8mm]}, thick, pcbmid},
  ctl/.style  ={-{Stealth[length=1.8mm]}, densely dashed, thick, muted},
  irq/.style  ={-{Stealth[length=2mm]}, densely dotted, line width=1.3pt, bad},
  num/.style  ={circle, fill=copper, text=white, font=\scriptsize\bfseries, inner sep=1pt,
                minimum size=0.46cm},
  ours/.style ={draw=bad, line width=1.2pt, densely dashed, rounded corners=4pt},
}

% ---------------- helpers ----------------
\newcommand{\rn}[1]{\texttt{#1}}
% \ffrow{x}{y}{n}{cell}{style}: n flip-flops in a row, bottom-left corner at (x,y)
\newcommand{\ffrow}[5]{\foreach \q in {1,...,#3}{\draw[#5] ({#1+(\q-1)*#4},{#2}) rectangle ++(#4,#4);}}
% \ffcol{x}{y}{n}{cell}{style}: n flip-flops in a column, top-left corner at (x,y)
\newcommand{\ffcol}[5]{\foreach \q in {1,...,#3}{\draw[#5] ({#1},{#2-\q*#4}) rectangle ++(#4,#4);}}
% \ffgrid{x}{ytop}{n}{per row}{cell}{style}: n flip-flops, filled row by row from the top-left
\newcommand{\ffgrid}[6]{\pgfmathtruncatemacro{\gl}{#3-1}%
  \foreach \b in {0,...,\gl}{\pgfmathtruncatemacro{\gc}{mod(\b,#4)}\pgfmathtruncatemacro{\gr}{div(\b,#4)}%
    \draw[#6] ({#1+\gc*#5},{#2-(\gr+1)*#5}) rectangle ++(#5,#5);}}

% every figure: same page size, title, footer
\newcommand{\canvas}{\useasboundingbox (-0.3,-0.2) rectangle (38.6,30.4);}
\newcommand{\figtitle}[3]{%
  \node[num, minimum size=0.9cm, font=\Large\bfseries] at (0.45,29.55) {#1};
  \node[anchor=west, font=\LARGE\bfseries, text=pcb] at (1.15,29.55) {#2};
  \node[anchor=west, font=\small, text=ink] at (0,28.7) {#3};}
\newcommand{\figfoot}[1]{\node[anchor=west, font=\scriptsize, text=muted] at (0,0.25) {#1};}
% \panel{x0}{y0}{x1}{y1}{title}
\newcommand{\panel}[5]{\draw[panel] (#1,#2) rectangle (#3,#4);
  \node[ptitle] at ({#1+0.3},{#4-0.25}) {#5};}

% ---------------- timing diagrams: 1 cycle = \cw cm, low at y, high at y+0.45 ----------------
\newcommand{\cw}{1.0}
\newcommand{\wclk}[3]{% x0 y ncycles
  \foreach \k in {1,...,#3}{\draw[pcb, thick] ({#1+(\k-1)*\cw},#2) -- ++(0,0.45) -- ++(0.5*\cw,0)
     -- ++(0,-0.45) -- ++(0.5*\cw,0);}}
\newcommand{\wsig}[3]{% x0 y {level per cycle, e.g. 0,1,1,0}
  \gdef\wprev{x}%
  \foreach \v [count=\k from 0] in {#3}{%
    \draw[pcb, thick] ({#1+\k*\cw},{#2+0.45*\v}) -- ++(\cw,0);
    \if x\wprev\else\ifnum\wprev=\v\else\draw[pcb, thick] ({#1+\k*\cw},#2) -- ++(0,0.45);\fi\fi
    \global\let\wprev\v}}
\newcommand{\wbus}[5]{% x0 y c0 c1 text: one value held from cycle c0 to cycle c1
  \draw[pcbmid, thick, fill=white] ({#1+(#3)*\cw},{#2+0.225}) -- ({#1+(#3)*\cw+0.08},{#2+0.45})
    -- ({#1+(#4)*\cw-0.08},{#2+0.45}) -- ({#1+(#4)*\cw},{#2+0.225}) -- ({#1+(#4)*\cw-0.08},#2)
    -- ({#1+(#3)*\cw+0.08},#2) -- cycle;
  \node[font=\tiny\ttfamily] at ({#1+((#3)+(#4))*\cw/2},{#2+0.225}) {#5};}
\newcommand{\widle}[4]{% x0 y c0 c1: bus idle (no value)
  \draw[muted, thick] ({#1+(#3)*\cw},{#2+0.225}) -- ({#1+(#4)*\cw},{#2+0.225});}
\newcommand{\wname}[3]{\node[anchor=east, font=\scriptsize\ttfamily] at ({#1-0.15},{#2+0.225}) {#3};}
.
\usepackage[T1]{fontenc}
\usepackage[scaled=0.95]{helvet}
\renewcommand{\familydefault}{\sfdefault}
\usepackage{booktabs}
\usepackage{amsmath}
\usepackage{amssymb}
\usetikzlibrary{arrows.meta,calc,positioning,fit,shapes.geometric,shapes.misc}

% ---------------- palette ----------------
\definecolor{pcb}{HTML}{0B2E33}
\definecolor{pcbmid}{HTML}{1F5C61}
\definecolor{tint}{HTML}{E6F0EF}
\definecolor{copper}{HTML}{B5651D}
\definecolor{coppert}{HTML}{F6E8DC}
\definecolor{ink}{HTML}{1F2A2E}
\definecolor{muted}{HTML}{5B6B70}
\definecolor{hardf}{HTML}{DCE6F2}
\definecolor{harde}{HTML}{2F4A7A}
\definecolor{good}{HTML}{2E7D32}
\definecolor{bad}{HTML}{B3261E}
% one colour per kind of block: border (d) / fill (f)
\definecolor{cpud}{HTML}{2B5C9E} \definecolor{cpuf}{HTML}{D6E4F5}
\definecolor{gpud}{HTML}{2E7D32} \definecolor{gpuf}{HTML}{D5EDD0}
\definecolor{npud}{HTML}{0B2E33} \definecolor{npuf}{HTML}{BFE0DB}
\definecolor{memd}{HTML}{5E3A8C} \definecolor{memf}{HTML}{E4DAF3}
\definecolor{busd}{HTML}{B5651D} \definecolor{busf}{HTML}{F6E8DC}
\definecolor{prd}{HTML}{546E7A}  \definecolor{prf}{HTML}{E3EAED}
\definecolor{secd}{HTML}{B3261E} \definecolor{secf}{HTML}{F7D9D6}
\definecolor{pwrd}{HTML}{8A6800} \definecolor{pwrf}{HTML}{FFF0B8}
\definecolor{medd}{HTML}{8E2463} \definecolor{medf}{HTML}{F5D6E6}
\definecolor{skd}{HTML}{2E7D32}  \definecolor{skf}{HTML}{BFE3B4}   % skew buffers, as in the register figure

\tikzset{
  blk/.style 2 args={draw=#1, fill=#2, rounded corners=4pt, thick, align=center,
                     font=\scriptsize, inner sep=3pt},
  cpu/.style={blk={cpud}{cpuf}},  gpu/.style={blk={gpud}{gpuf}},
  npu/.style={blk={npud}{npuf}},  mem/.style={blk={memd}{memf}},
  bus/.style={blk={busd}{busf}},  per/.style={blk={prd}{prf}},
  sec/.style={blk={secd}{secf}},  pwr/.style={blk={pwrd}{pwrf}},
  med/.style={blk={medd}{medf}},
  comb/.style ={draw=copper, fill=coppert, rounded corners=2pt, thick, align=center,
                font=\scriptsize, inner sep=2pt},                       % combinational logic
  hard/.style ={draw=harde, fill=hardf, thick, dashed, align=center,
                font=\scriptsize, inner sep=2pt},                       % hard macro / off-chip
  sram/.style ={draw=memd, fill=memf, thick, align=center, font=\scriptsize,
                inner sep=2pt},                                         % SRAM array, not flip-flops
  ff/.style   ={draw=#1, fill=#1!30, line width=0.35pt},                % one flip-flop square
  ffgone/.style={draw=muted, fill=white, densely dashed, line width=0.35pt},
  panel/.style={draw=muted, rounded corners=6pt, line width=0.9pt},
  ptitle/.style={anchor=north west, font=\normalsize\bfseries, text=ink, inner sep=0pt},
  note/.style ={font=\scriptsize, text=ink, align=left, inner sep=0pt},
  lab/.style  ={font=\tiny, text=muted, inner sep=1pt},
  axi/.style  ={{Stealth[length=2.6mm]}-{Stealth[length=2.6mm]}, line width=2.2pt, copper},
  axi1/.style ={-{Stealth[length=2.6mm]}, line width=2.2pt, copper},
  apb/.style  ={{Stealth[length=1.8mm]}-{Stealth[length=1.8mm]}, line width=1pt, copper},
  apb1/.style ={-{Stealth[length=1.8mm]}, line width=1pt, copper},
  dat/.style  ={-{Stealth[length=1.8mm]}, thick, pcbmid},
  ctl/.style  ={-{Stealth[length=1.8mm]}, densely dashed, thick, muted},
  irq/.style  ={-{Stealth[length=2mm]}, densely dotted, line width=1.3pt, bad},
  num/.style  ={circle, fill=copper, text=white, font=\scriptsize\bfseries, inner sep=1pt,
                minimum size=0.46cm},
  ours/.style ={draw=bad, line width=1.2pt, densely dashed, rounded corners=4pt},
}

% ---------------- helpers ----------------
\newcommand{\rn}[1]{\texttt{#1}}
% \ffrow{x}{y}{n}{cell}{style}: n flip-flops in a row, bottom-left corner at (x,y)
\newcommand{\ffrow}[5]{\foreach \q in {1,...,#3}{\draw[#5] ({#1+(\q-1)*#4},{#2}) rectangle ++(#4,#4);}}
% \ffcol{x}{y}{n}{cell}{style}: n flip-flops in a column, top-left corner at (x,y)
\newcommand{\ffcol}[5]{\foreach \q in {1,...,#3}{\draw[#5] ({#1},{#2-\q*#4}) rectangle ++(#4,#4);}}
% \ffgrid{x}{ytop}{n}{per row}{cell}{style}: n flip-flops, filled row by row from the top-left
\newcommand{\ffgrid}[6]{\pgfmathtruncatemacro{\gl}{#3-1}%
  \foreach \b in {0,...,\gl}{\pgfmathtruncatemacro{\gc}{mod(\b,#4)}\pgfmathtruncatemacro{\gr}{div(\b,#4)}%
    \draw[#6] ({#1+\gc*#5},{#2-(\gr+1)*#5}) rectangle ++(#5,#5);}}

% every figure: same page size, title, footer
\newcommand{\canvas}{\useasboundingbox (-0.3,-0.2) rectangle (38.6,30.4);}
\newcommand{\figtitle}[3]{%
  \node[num, minimum size=0.9cm, font=\Large\bfseries] at (0.45,29.55) {#1};
  \node[anchor=west, font=\LARGE\bfseries, text=pcb] at (1.15,29.55) {#2};
  \node[anchor=west, font=\small, text=ink] at (0,28.7) {#3};}
\newcommand{\figfoot}[1]{\node[anchor=west, font=\scriptsize, text=muted] at (0,0.25) {#1};}
% \panel{x0}{y0}{x1}{y1}{title}
\newcommand{\panel}[5]{\draw[panel] (#1,#2) rectangle (#3,#4);
  \node[ptitle] at ({#1+0.3},{#4-0.25}) {#5};}

% ---------------- timing diagrams: 1 cycle = \cw cm, low at y, high at y+0.45 ----------------
\newcommand{\cw}{1.0}
\newcommand{\wclk}[3]{% x0 y ncycles
  \foreach \k in {1,...,#3}{\draw[pcb, thick] ({#1+(\k-1)*\cw},#2) -- ++(0,0.45) -- ++(0.5*\cw,0)
     -- ++(0,-0.45) -- ++(0.5*\cw,0);}}
\newcommand{\wsig}[3]{% x0 y {level per cycle, e.g. 0,1,1,0}
  \gdef\wprev{x}%
  \foreach \v [count=\k from 0] in {#3}{%
    \draw[pcb, thick] ({#1+\k*\cw},{#2+0.45*\v}) -- ++(\cw,0);
    \if x\wprev\else\ifnum\wprev=\v\else\draw[pcb, thick] ({#1+\k*\cw},#2) -- ++(0,0.45);\fi\fi
    \global\let\wprev\v}}
\newcommand{\wbus}[5]{% x0 y c0 c1 text: one value held from cycle c0 to cycle c1
  \draw[pcbmid, thick, fill=white] ({#1+(#3)*\cw},{#2+0.225}) -- ({#1+(#3)*\cw+0.08},{#2+0.45})
    -- ({#1+(#4)*\cw-0.08},{#2+0.45}) -- ({#1+(#4)*\cw},{#2+0.225}) -- ({#1+(#4)*\cw-0.08},#2)
    -- ({#1+(#3)*\cw+0.08},#2) -- cycle;
  \node[font=\tiny\ttfamily] at ({#1+((#3)+(#4))*\cw/2},{#2+0.225}) {#5};}
\newcommand{\widle}[4]{% x0 y c0 c1: bus idle (no value)
  \draw[muted, thick] ({#1+(#3)*\cw},{#2+0.225}) -- ({#1+(#4)*\cw},{#2+0.225});}
\newcommand{\wname}[3]{\node[anchor=east, font=\scriptsize\ttfamily] at ({#1-0.15},{#2+0.225}) {#3};}

\begin{document}
\begin{tikzpicture}[line join=round]
\canvas
\figtitle{1}{The whole SoC: every block on one chip}
  {A typical phone / edge-AI system-on-chip. Orange number = the figure that zooms into that block.
   Thick copper = AXI data bus, thin copper = APB / control bus, red dotted = interrupt wire.}

% ---------------- SoC boundary ----------------
\draw[pcb, line width=1.6pt, rounded corners=8pt] (0,8.6) rectangle (38.3,28.1);
\node[anchor=north west, font=\small\bfseries, text=pcb] at (0.3,27.95) {SoC = one silicon die};
\node[anchor=north east, font=\small, text=muted] at (38.0,27.95) {all sizes are typical examples};

% =====================================================================
% row 1: compute engines
% =====================================================================
% ---- CPU cluster ----
\draw[cpu] (0.4,19.6) rectangle (10.6,27.3);
\node[anchor=north west, font=\normalsize\bfseries, text=cpud] at (0.6,27.15) {CPU cluster};
\node[num] at (8.9,26.95) {2}; \node[num] at (9.5,26.95) {3};
\foreach \k in {0,1}{
  \node[cpu, fill=white, minimum width=2.3cm, minimum height=2.7cm, text width=2.1cm, anchor=south west]
    at ({0.8+\k*2.5},23.55) {\textbf{big core \k}\\out-of-order\\[2pt]L1-I 64 KB\\L1-D 64 KB\\L2 1 MB};}
\foreach \k in {0,...,3}{
  \node[cpu, fill=white, minimum width=1.05cm, minimum height=2.7cm, text width=0.9cm, anchor=south west,
        font=\tiny] at ({5.85+\k*1.13},23.55) {\textbf{little}\\\textbf{core \k}\\in-order\\[2pt]L1\\32+32\\KB};}
\node[sram, minimum width=9.5cm, minimum height=1.9cm, text width=9.2cm, anchor=south west]
  at (0.8,21.3) {\textbf{shared L3 cache 8 MB} $\cdot$ snoop filter $\cdot$ cluster control unit\\
                 keeps the 6 cores' caches coherent (same value everywhere)};
\node[bus, minimum width=9.5cm, minimum height=0.9cm, anchor=south west] at (0.8,20.0)
  {coherent port to the NoC (CHI / ACE)};

% ---- GPU ----
\draw[gpu] (11.0,19.6) rectangle (19.4,27.3);
\node[anchor=north west, font=\normalsize\bfseries, text=gpud] at (11.2,27.15) {GPU};
\node[num] at (18.9,26.95) {4};
\node[gpu, fill=white, minimum width=7.8cm, minimum height=0.95cm, anchor=south west] at (11.3,25.45)
  {command processor $\cdot$ work dispatcher $\cdot$ GPU MMU};
\foreach \r in {0,1}{\foreach \c in {0,...,3}{
  \node[gpu, fill=white, minimum width=1.8cm, minimum height=1.25cm, anchor=south west, font=\tiny]
    at ({11.3+\c*2.0},{22.75+\r*1.38}) {\textbf{shader core}\\64 lanes};}}
\node[comb, minimum width=7.8cm, minimum height=0.75cm, anchor=south west] at (11.3,21.75)
  {fixed function: tiler $\cdot$ rasterizer $\cdot$ texture $\cdot$ blend};
\node[sram, minimum width=7.8cm, minimum height=1.25cm, anchor=south west] at (11.3,20.0)
  {\textbf{GPU L2 cache 2 MB}\\shared by all shader cores};

% ---- NPU ----
\draw[npu] (19.8,19.6) rectangle (29.0,27.3);
\node[anchor=north west, font=\normalsize\bfseries, text=npud] at (20.0,27.15) {NPU (AI accelerator)};
\node[num] at (28.5,26.95) {5};
\node[npu, fill=white, minimum width=2.6cm, minimum height=1.05cm, anchor=south west, font=\tiny]
  at (20.1,25.35) {\textbf{control regs}\\AXI-Lite + IRQ};
\node[npu, fill=white, minimum width=2.7cm, minimum height=1.05cm, anchor=south west, font=\tiny]
  at (22.9,25.35) {\textbf{command}\\\textbf{processor}};
\node[npu, fill=white, minimum width=2.7cm, minimum height=1.05cm, anchor=south west, font=\tiny]
  at (25.8,25.35) {\textbf{DMA}\\AXI master};
\foreach \nm [count=\k from 0] in {weights, inputs, outputs}{
  \node[sram, minimum width=2.75cm, minimum height=0.85cm, anchor=south west, font=\tiny]
    at ({20.1+\k*2.88},24.3) {\textbf{SRAM: \nm}\\ping $|$ pong};}
\foreach \r in {0,...,7}{\foreach \c in {0,...,7}{
  \draw[ff=npud] ({20.75+\c*0.42},{20.35+\r*0.42}) rectangle ++(0.34,0.34);}}
\draw[dat] (20.15,22.05) -- (20.7,22.05);
\draw[dat] (22.4,24.25) -- (22.4,23.75);
\node[lab, text=npud, font=\tiny\bfseries, anchor=south] at (22.4,19.95) {MAC array 32$\times$32 int8 (drawn 8$\times$8)};
\draw[ours] (20.6,20.2) rectangle (24.25,23.8);
\node[npu, fill=white, minimum width=4.1cm, minimum height=0.6cm, anchor=south west, font=\tiny] at (24.6,23.25)
  {accumulators int32 (registers)};
\node[comb, minimum width=4.1cm, minimum height=0.6cm, anchor=south west, font=\tiny] at (24.6,22.5)
  {+ bias $\rightarrow$ requantize to int8};
\node[comb, minimum width=4.1cm, minimum height=0.6cm, anchor=south west, font=\tiny] at (24.6,21.75)
  {activation (ReLU) $\rightarrow$ pooling};
\node[anchor=north west, font=\tiny, text=bad, text width=4.1cm, align=left] at (24.6,21.55)
  {\textbf{red dashed} = what your \rn{mm\_core} already is:\\a 4$\times$4 MAC array + control regs\\+ small RAM buffers};

% ---- media ----
\draw[med] (29.4,19.6) rectangle (37.9,27.3);
\node[anchor=north west, font=\normalsize\bfseries, text=medd] at (29.6,27.15) {Media and sensing};
\foreach \nm/\tx [count=\k from 0] in {
    ISP (camera)/{demosaic, denoise, HDR, 3A},
    DSP/{audio, sensors, always-on voice},
    video codec/{H.265 / AV1 encode + decode},
    display controller/{layers, scaling, MIPI DSI out}}{
  \node[med, fill=white, minimum width=7.9cm, minimum height=1.45cm, text width=7.6cm, anchor=south west]
    at (29.7,{24.85-\k*1.62}) {\textbf{\nm}\\\tx};}

% =====================================================================
% row 2: interconnect
% =====================================================================
\draw[bus] (0.4,16.3) rectangle (37.9,18.95);
\node[anchor=north west, font=\normalsize\bfseries, text=busd] at (0.6,18.85) {On-chip interconnect (network-on-chip)};
\node[num] at (9.6,18.55) {6};
\foreach \x in {3.0,9.0,15.2,21.6,27.6,32.4}{
  \node[draw=busd, fill=white, thick, minimum size=0.8cm, font=\tiny\bfseries] (r\x) at (\x,17.15) {R};}
\draw[busd, line width=2pt] (3.4,17.15) -- (8.6,17.15) (9.4,17.15) -- (14.8,17.15) (15.6,17.15) -- (21.2,17.15)
  (22.0,17.15) -- (27.2,17.15) (28.0,17.15) -- (32.0,17.15) (32.8,17.15) -- (33.9,17.15);
\node[font=\tiny, text=busd] at (6.0,16.6) {R = router: buffers + arbiter + switch};
\node[sram, minimum width=3.8cm, minimum height=1.6cm, anchor=south west, font=\tiny] at (33.9,16.45)
  {\textbf{system level cache}\\\textbf{8 MB}\\shared by CPU, GPU, NPU};
\node[font=\tiny, text=busd, align=center] at (18.4,16.6) {address decode $\cdot$ arbitration $\cdot$ QoS (who goes first)};

% compute engines -> NoC
\foreach \x in {5.5,15.2,24.4,33.6}{\draw[axi] (\x,19.6) -- (\x,18.95);}

% =====================================================================
% row 3: system blocks
% =====================================================================
% ---- memory controller ----
\draw[mem] (0.4,9.0) rectangle (8.4,15.6);
\node[anchor=north west, font=\normalsize\bfseries, text=memd] at (0.6,15.45) {DRAM controller};
\node[num] at (7.9,15.2) {7};
\node[mem, fill=white, minimum width=7.4cm, minimum height=1.35cm, text width=7.1cm, anchor=south west] at (0.7,12.85)
  {\textbf{front-end}: AXI ports, QoS, read / write queues, address mapping};
\node[mem, fill=white, minimum width=7.4cm, minimum height=1.35cm, text width=7.1cm, anchor=south west] at (0.7,11.3)
  {\textbf{scheduler}: bank state machines, timing counters, refresh timer};
\node[hard, minimum width=7.4cm, minimum height=1.35cm, text width=7.1cm, anchor=south west] at (0.7,9.3)
  {\textbf{PHY} (analog): serializers, delay lines, I/O pads, link training};

% ---- DMA + GIC ----
\draw[per] (8.8,9.0) rectangle (15.6,15.6);
\node[anchor=north west, font=\normalsize\bfseries, text=prd] at (9.0,15.45) {DMA + interrupts};
\node[num] at (15.1,15.2) {8};
\node[per, fill=white, minimum width=6.2cm, minimum height=2.3cm, text width=5.9cm, anchor=south west] at (9.1,12.2)
  {\textbf{system DMA}: 8 channels\\copies memory $\rightarrow$ memory or memory $\leftrightarrow$ device without the CPU, following descriptors};
\node[per, fill=white, minimum width=6.2cm, minimum height=2.6cm, text width=5.9cm, anchor=south west] at (9.1,9.3)
  {\textbf{GIC} (interrupt controller)\\hundreds of IRQ wires in $\rightarrow$ enable, priority, target $\rightarrow$ one IRQ to the chosen CPU core};

% ---- clocks, reset, power ----
\draw[pwr] (16.0,9.0) rectangle (23.2,15.6);
\node[anchor=north west, font=\normalsize\bfseries, text=pwrd] at (16.2,15.45) {Clocks, reset, power};
\node[num] at (22.7,15.2) {9};
\foreach \nm/\tx [count=\k from 0] in {
    PLLs/{multiply 24 MHz up to 3 GHz},
    clock gating/{stop clocks of idle blocks},
    reset controller/{assert async, release in sync},
    PMU + DVFS/{power domains on/off, voltage + frequency}}{
  \node[pwr, fill=white, minimum width=6.6cm, minimum height=1.25cm, text width=6.3cm, anchor=south west]
    at (16.3,{13.0-\k*1.35}) {\textbf{\nm}: \tx};}

% ---- security, boot, debug ----
\draw[sec] (23.6,9.0) rectangle (30.6,15.6);
\node[anchor=north west, font=\normalsize\bfseries, text=secd] at (23.8,15.45) {Boot, security, debug};
\node[num] at (30.1,15.2) {10};
\foreach \nm/\tx [count=\k from 0] in {
    boot ROM + SRAM/{first code after reset},
    OTP fuses/{root-key hash, lock bits},
    crypto + TRNG/{AES, SHA, RSA/ECC, random},
    TrustZone + debug/{secure world, JTAG, trace}}{
  \node[sec, fill=white, minimum width=6.4cm, minimum height=1.25cm, text width=6.1cm, anchor=south west]
    at (23.9,{13.0-\k*1.35}) {\textbf{\nm}: \tx};}

% ---- peripherals ----
\draw[per] (31.0,9.0) rectangle (37.9,15.6);
\node[anchor=north west, font=\normalsize\bfseries, text=prd] at (31.2,15.45) {Peripherals};
\node[num] at (37.4,15.2) {11};
\node[bus, minimum width=6.3cm, minimum height=0.8cm, anchor=south west] at (31.3,13.6) {AXI $\rightarrow$ APB bridge};
\foreach \nm [count=\k from 0] in {UART, SPI, I2C, GPIO, timers, watchdog}{
  \pgfmathtruncatemacro{\cc}{mod(\k,3)}\pgfmathtruncatemacro{\rr}{div(\k,3)}
  \node[per, fill=white, minimum width=1.95cm, minimum height=0.8cm, anchor=south west, font=\tiny]
    at ({31.3+\cc*2.12},{12.55-\rr*0.95}) {\nm};}
\node[per, fill=white, minimum width=6.3cm, minimum height=1.9cm, text width=6.0cm, anchor=south west] at (31.3,9.3)
  {\textbf{high-speed I/O} (own PHYs):\\USB 3 $\cdot$ PCIe $\cdot$ UFS storage\\MIPI CSI-2 camera $\cdot$ MIPI DSI display};

% NoC -> row 3
\draw[axi] (4.4,16.3) -- (4.4,15.6);
\draw[axi] (12.2,16.3) -- (12.2,15.6);
\draw[apb] (19.6,16.3) -- (19.6,15.6);
\draw[apb] (27.1,16.3) -- (27.1,15.6);
\draw[axi] (34.4,16.3) -- (34.4,15.6);

% interrupt wires (dedicated wires, not through the NoC)
\draw[irq] (20.3,19.6) -- (20.3,15.95) -- (14.6,15.95) -- (14.6,15.6);
\draw[irq] (10.2,15.6) -- (10.2,19.6);
\node[font=\tiny\bfseries, text=bad, anchor=west, fill=white, inner sep=1pt] at (15.0,15.95) {IRQ wire from NPU};
\node[font=\tiny\bfseries, text=bad, anchor=east, fill=busf, inner sep=1pt, rotate=90] at (9.95,17.9) {IRQ to CPU};

% =====================================================================
% off-chip parts
% =====================================================================
\node[anchor=west, font=\small\bfseries, text=harde] at (0,8.15) {off-chip (on the board)};
\node[hard, minimum width=8.0cm, minimum height=1.9cm, text width=7.7cm, anchor=south west] at (0.4,5.2)
  {\textbf{LPDDR5 DRAM} 8--16 GB\\4 channels $\times$ 16 bit, 6400 MT/s $\approx$ 51 GB/s};
\node[hard, minimum width=3.3cm, minimum height=1.9cm, text width=3.1cm, anchor=south west] at (16.0,5.2)
  {\textbf{PMIC}\\makes 10+ supply rails};
\node[hard, minimum width=3.4cm, minimum height=1.9cm, text width=3.2cm, anchor=south west] at (19.8,5.2)
  {\textbf{crystal}\\24 MHz reference};
\node[hard, minimum width=3.3cm, minimum height=1.9cm, text width=3.1cm, anchor=south west] at (23.6,5.2)
  {\textbf{JTAG probe}\\lab / debug only};
\node[hard, minimum width=3.3cm, minimum height=1.9cm, text width=3.1cm, anchor=south west] at (27.3,5.2)
  {\textbf{UFS flash}\\128 GB -- 1 TB};
\node[hard, minimum width=3.3cm, minimum height=1.9cm, text width=3.1cm, anchor=south west] at (31.0,5.2)
  {\textbf{camera}\\image sensor};
\node[hard, minimum width=3.3cm, minimum height=1.9cm, text width=3.1cm, anchor=south west] at (34.6,5.2)
  {\textbf{display}\\panel};
\draw[axi] (6.8,9.0) -- (6.8,7.1);
\draw[apb] (17.65,9.0) -- (17.65,7.1);
\draw[apb1] (21.5,7.1) -- (21.5,9.0);
\draw[apb] (25.25,9.0) -- (25.25,7.1);
\draw[apb] (31.6,9.0) -- (31.6,8.4) -- (28.95,8.4) -- (28.95,7.1);
\draw[apb] (32.65,9.0) -- (32.65,7.1);
\draw[apb] (36.25,9.0) -- (36.25,7.1);

% =====================================================================
% bottom panels
% =====================================================================
\panel{0}{0.7}{12.5}{4.75}{How to read this map}
\node[note, font=\small, anchor=north west, text width=11.9cm] at (0.3,4.05) {%
  Colour = kind of block: \textcolor{cpud}{\textbf{CPU}}, \textcolor{gpud}{\textbf{GPU}}, \textcolor{npud}{\textbf{NPU}},
  \textcolor{memd}{\textbf{memory}}, \textcolor{busd}{\textbf{interconnect}}, \textcolor{prd}{\textbf{peripherals / DMA}},
  \textcolor{secd}{\textbf{security}}, \textcolor{pwrd}{\textbf{clocks / power}}, \textcolor{medd}{\textbf{media}}.\\[2pt]
  Purple box = SRAM array. Peach = combinational logic. Dashed blue = analog or off-chip.\\[2pt]
  Every block is built from the same three things as your \rn{mm\_core}: flip-flops, combinational logic and SRAM.};
\panel{12.8}{0.7}{25.4}{4.75}{One AI job in 6 steps (Figure 8 has every detail)}
\node[note, font=\small, anchor=north west, text width=12.0cm] at (13.1,4.05) {%
  \textbf{1} CPU puts the input and the weights in DRAM.\quad
  \textbf{2} CPU writes the NPU's START register (AXI-Lite).\quad
  \textbf{3} NPU DMA reads the data from DRAM through the NoC into its SRAM.\quad
  \textbf{4} MAC array computes, post-processing makes int8 outputs.\quad
  \textbf{5} DMA writes the result to DRAM.\quad
  \textbf{6} NPU raises its IRQ wire $\rightarrow$ GIC $\rightarrow$ CPU reads the result.};
\panel{25.7}{0.7}{38.3}{4.75}{Where your project fits}
\node[note, font=\small, anchor=north west, text width=12.0cm] at (26.0,4.05) {%
  \rn{mm\_core} = the NPU's \textbf{MAC array} (4$\times$4 instead of 32$\times$32) + \textbf{control registers}
  (\rn{mm\_regs}) + small \textbf{RAM buffers} (8 M9K).\\[2pt]
  UART + bridge = the bus. \rn{host.py} = the driver on the CPU. Polling STATUS = instead of the IRQ.\\[2pt]
  Not in your design yet: DMA, command processor, post-processing, big SRAM.};

\figfoot{Sizes and speeds are typical values for a recent phone / edge-AI SoC, not one specific chip. Figures 2--13 zoom into each numbered block.}
\end{tikzpicture}
\end{document}
